% Part: sets-functions-relations
如果一个集合与自然数集 $\Nat$ 之间存在双射，则称该集合是\emph{可数无穷的}（\emph{denumerable}）。
等价地，可数无穷集可以看作一个能按第一个元素、第二个元素……的方式列出其所有元素的集合。
由于 $\Nat$ 是最初的或者说典型的无穷集，因此使用“denumerable”这个词。
有时人们也使用“countably infinite”这个术语。

例如，自然数集合 $\Nat$ 显然是可数无穷的；
整数集 $\Int$ 和有理数集 $\Rat$ 也是可数无穷的。
（$\Int$ 的可数性很容易证明；至于 $\Rat$，可数性的证明需要更多论证，我们会在后面给出。）

但并非所有无穷集都是可数无穷的。
有些集合，比如实数集 $\Real$，是\emph{不可数的}。
所有由 $0$ 和 $1$ 构成的无限序列的集合 $\mathbf{2}^{\Nat}$ 就不是可数无穷的，这一点可以通过\emph{对角线论证}法来证明。
而 $\Real$ 是不可数的，因为 $\Real$ 的一个子集与 $\mathbf{2}^{\Nat}$ 可建立双射（即，每个无限 $0$-$1$ 序列都可以被看作是某个实数的二进制展开）。
我们将在后面给出完整的证明。
% Chapter: functions
% Section: basics

\documentclass[../../../include/open-logic-section]{subfiles}

\begin{document}

\olfileid{sfr}{fun}{bas}
\olsection{基本概念}

\begin{explain}
\emph{函数}是一种映射，它将给定集合中的每个元素映射到某个（另一个）给定集合中的特定元素。例如，加~$1$ 的运算定义了一个函数：每个数~$n$ 被映射到唯一的数~$n+1$。

更一般地，函数可以以有序对、三元组等为输入，并返回某种输出。我们从基础算术中就已经熟悉许多函数。例如，加法和乘法都是函数。它们接受两个数作为输入，返回第三个数。

在这种抽象的数学意义下，函数就是一个\emph{黑箱}：重要的仅仅是哪个输入与哪个输出配对，而非计算输出的方法。
\end{explain}

\begin{defn}[函数]
\emph{函数} $f \colon A \to B$ 指的是将~$A$ 的每个元素映射到~$B$ 的一个元素的映射。

我们称 $A$ 为~$f$ 的\emph{定义域}，$B$ 为~$f$ 的\emph{陪域}。$A$ 的元素称为~$f$ 的输入或\emph{自变元}，而由~$f$ 与自变元~$x$ 相配对的~$B$ 中的元素称为~$f$ 在自变元~$x$ 处的\emph{值}，记作~$f(x)$。

$f$ 的\emph{值域} $\ran{f}$ 是陪域的一个子集，由 $f$ 在某自变元处取到的所有值组成；$\ran{f} =
\Setabs{f(x)}{x \in A}$。
\end{defn}

\olref{fig:function} 中的图示可能有助于思考函数。左侧的椭圆代表函数的\emph{定义域}；右侧的椭圆代表函数的\emph{陪域}；箭头从定义域中的一个\emph{自变元}指向陪域中对应的\emph{值}。

\begin{figure}
  \olasset{assets/diagrams/function.tikz}
  \caption{函数是将一个集合的每个元素映射到另一个集合中某个元素的映射。箭头从定义域中的一个自变元指向陪域中对应的值。}
  \ollabel{fig:function}
\end{figure}

\begin{ex}
乘法以自然数对作为输入，并将它们映射到自然数作为输出，因此它是从 $\Nat \times \Nat$（定义域）到 $\Nat$（陪域）的映射。事实上，它的值域也是 $\Nat$，因为每个 $n \in \Nat$ 都可以表示为 $n \times 1$。
\end{ex}

\begin{ex}
乘法是一个函数，因为它将每个输入（每个自然数对）与唯一的一个输出配对：$\times \colon \Nat^2 \to \Nat$。相比之下，将一个自然数~$n$ 映射到一个满足 $x^2 = n$ 的实数~$x$ 则不是函数，因为每个正整数~$n$ 都有两个平方根：$\sqrt{n}$ 和~$-\sqrt{n}$。我们可以通过只返回正平方根使其成为函数：$\sqrt{\phantom{X}} \colon
\Nat \to \Real$。
\end{ex}

\begin{ex}
将班级里每个学生与其最终成绩配对的关系是一个函数——没有学生能在同一门课中获得两个不同的最终成绩。将班级里每个学生与其父母配对的关系则不是函数：学生可能有零个、两个或更多的父母。
\end{ex}

\begin{explain}
我们可以通过某种精确的方式，为每个可能的自变元指定函数的值，以此来定义函数。不同的定义方式包括：给出一个公式，描述一种计算值的方法，或者列出每个自变元对应的值。然而，无论函数如何定义，我们都必须确保为每个自变元指定一个，且仅有一个值。
\end{explain}


\begin{ex}
定义 $f \colon \Nat \to \Nat$ 使得 $f(x) = x+1$。这个定义将 $f$ 确定为一个以自然数为输入并以自然数为输出的函数。它告诉我们，给定一个自然数~$x$，$f$ 将输出它的后继~$x+1$。
在这种情况下，陪域 $\Nat$ 并不是~$f$ 的值域，因为自然数~$0$ 不是任何自然数的后继。$f$ 的值域是所有正整数的集合 $\Int^{+}$。
\end{ex}

\begin{ex}\ollabel{examplefunext}
定义 $g \colon \Nat \to \Nat$ 使得 $g(x) = x+2-1$。这告诉我们 $g$ 是一个以自然数为输入并输出自然数的函数。给定一个自然数~$n$，$g$ 将输出~$x$ 的后继的后继的前驱，即 $x+1$。
\end{ex}

\begin{explain}
我们刚刚考察了两个定义不同的函数 $f$ 和 $g$。然而，它们是\emph{同一个函数}。毕竟，对于任何自然数~$n$，我们有 $f(n) = n+1 = n+2-1 =
g(n)$。换句话说：我们对 $f$ 和~$g$ 的定义通过不同的方程指定了相同的映射。隐含地，我们实际上是依赖于函数的外延性原理，
\[
  \text{若 }\forall x\, f(x) = g(x)\text{, 则 }f = g
\]
前提是 $f$ 和~$g$ 具有相同的定义域和陪域。
\end{explain}

\begin{ex}
我们也可以用分情况的方式来定义函数。例如，我们可以如下定义 $h \colon \Nat \to \Nat$：
\[
h(x) =
\begin{cases}
  \frac{x}{2} & \text{若 $x$ 是偶数} \\
  \frac{x+1}{2} & \text{若 $x$ 是奇数。}
\end{cases}
\]
由于每个自然数要么是偶数要么是奇数，因此这个函数的输出总是一个自然数。只需记住，如果你用分情况的方式定义一个函数，那么每一个可能的输入都必须恰好落入一种情况。在某些情况下，这需要证明这些情况是穷尽且互斥的。
\end{ex}

\end{document}